\subsection{57. Can PINNs Beat the Finite Element Method? (Grossmann et al.\ IMA J.~Numer.~Anal.\ 2023)}
\textbf{arXiv:} \texttt{2302.04107}\quad\textbf{Authors:} Grossmann, Komorowska, Latz, Sch\"onlieb (2023)\quad\textbf{Domain:} Numerical PDE / ML\\
\scorebadge{Coverage}{8}\quad\scorebadge{Agreement}{9}\quad\textbf{Total:} 17/20\par\smallskip
\noindent\textbf{Coverage rationale.} 5 of 6 benchmark problems tested (83\%): 1D/2D/3D Poisson, 1D Allen--Cahn, and 1D Schr\"odinger. The 2D Schr\"odinger problem was not attempted (requires FEniCS for 2D ground truth). Two previously data-blocked problems (3D Poisson and 1D Schr\"odinger) were unblocked by self-generating ground truth: analytical solution for 3D Poisson, split-step Fourier (N=8192, $dt$=10$^{-5}$) for 1D Schr\"odinger. 11/12 testable claims evaluated.\par
\noindent\textbf{Agreement rationale.} 11/12 claims verified, 1 partial. FEM dominates PINNs on every tested problem: accuracy advantage 8.8$\times$--7900$\times$ in relative $L^2$ error; speed advantage 1.5$\times$--3900$\times$ in wall-clock time. Specific matches: 1D Poisson FEM achieves $3.80 \times 10^{-8}$ in 0.004\,s vs best PINN $2.02 \times 10^{-6}$ in 14.9\,s. 3D Poisson FEM achieves $1.68 \times 10^{-3}$ in 4.3\,s vs best PINN $1.47 \times 10^{-2}$ in 57\,s. PINN accuracy plateaus with increasing parameters, consistent with the paper.\par\smallskip
\noindent\textbf{Key findings:}
\begin{itemize}[leftmargin=1.4em,itemsep=0.2em,topsep=0.2em]
\item FEM beats PINNs on \emph{every} tested problem---accuracy and speed
\item 3D Poisson (NEW---previously data-blocked): FEM 8.8$\times$ more accurate, 13$\times$ faster. Width matters more than depth for PINNs; [60,60,1] best architecture.
\item 1D Schr\"odinger (NEW---previously data-blocked): FEM 43$\times$ more accurate, 18$\times$ faster. PINNs fail to capture soliton collision dynamics ($\sim$20\% rel.~error).
\item PINN non-monotonic scaling confirmed: larger networks sometimes \emph{worse} (1D Schr\"odinger [100,100,100,2] worse than [20,20,20,2])
\end{itemize}
\noindent\textbf{Verdict: REPLICATED.} Paper's central thesis---FEM consistently outperforms PINNs---is independently confirmed across 5 PDE benchmarks.
\noindent\textbf{Follow-on research questions:}
\begin{enumerate}[leftmargin=1.6em,itemsep=0.2em,topsep=0.2em,label=Q\arabic*.]
\item Does the FEM advantage persist for higher-order elements (CG2, CG3)? Repeat the comparison with $p$-refinement and measure whether PINNs become competitive at high polynomial order.
\item Can adaptive collocation point strategies (e.g., residual-based, NTK-weighted) close the PINN--FEM gap on the 1D Schr\"odinger benchmark?
\item How do Fourier Neural Operators (FNO) compare against both FEM and PINNs on the same 6-problem benchmark suite?
\item At what problem complexity (number of spatial dimensions, nonlinearity degree, irregular geometry) do PINNs begin to outperform FEM? Systematically increase dimension from 3D to 10D for Poisson.
\item Train PINNs with modern techniques (learning rate warmup, Sobolev training, causal weighting) and re-evaluate---does the $\sim$2 OOM accuracy gap close?
\end{enumerate}

